% 警告修复1：使用标准 article 类，去掉无效选项 ctexart 和 openany
%   ctexart 是独立文档类，不是 article 的选项；openany 是 book/report 专用
% 警告修复2：不手动加载 xeCJK 和 fontspec，也不手动设置字体
%   原因：\usepackage{ctex} 内部已自动处理 Windows 字体（SimSun/YaHei）
%         手动再次调用 \setCJKmainfont 会触发 xeCJK "Redefining CJKfamily" 警告
\documentclass[UTF8,a4paper,11pt]{article}
\usepackage{ctex}

\usepackage{mathtools}
\usepackage{amsmath}
\usepackage{amsfonts}
\usepackage{amssymb}
\usepackage{amsthm}
\usepackage[T1]{fontenc}
\usepackage{indentfirst}

\usepackage{graphicx}
\usepackage{geometry}
\usepackage{latexsym}
\usepackage{fancyhdr}
\usepackage{titlesec}
\usepackage{setspace}
\usepackage{enumerate}
\usepackage{enumitem}
\usepackage{multicol}
\usepackage{url}
\usepackage{ulem}
\usepackage{bm}
\usepackage{hyperref}
\usepackage{stfloats}

\setlength{\parindent}{2em}
\linespread{1.2}
\usepackage{listings}
\usepackage{xcolor}
\usepackage{algorithm}
\usepackage{algorithmicx}
\usepackage{algpseudocode}
\usepackage{mdframed}
\usepackage{booktabs}
\usepackage{multirow}  % 支持表格中跨行合并 \multirow{行数}{宽度}{内容}
\usepackage{float}     % 支持浮动体选项 [H]（强制在此处放置）
\usepackage{makecell}  % 支持 \makecell 在单元格内换行

\everymath{\displaystyle}

\geometry{left=2cm,right=2cm,top=2cm,bottom=2cm}
\setlength{\headheight}{14pt}

\pagestyle{fancy}
\fancyfoot[C]{}
% 警告修复3：单面模式只用 O（奇数页）选项，去掉无效的 E（偶数页）选项
\fancyhead[RO]{\thepage}
\fancyhead[LO]{\leftmark\quad 学号：241098117，姓名：张城宗}

\titleformat{\section}{\bfseries\Large}{$\S$\,\thesection\,}{1em}{}
\setenumerate[1]{itemsep=0pt,partopsep=0pt,parsep=\parskip,topsep=0pt}
\setitemize[1]{itemsep=0pt,partopsep=0pt,parsep=\parskip,topsep=0pt}

\newcommand*{\rmk}{\textbf{注：}}
\newcommand\e{\leftarrow}

% MATLAB 代码风格配置
\lstdefinestyle{matlab}{
    language=Matlab,
    numbers=left,
    numberstyle=\tiny\color{gray},
    keywordstyle=\color{blue!80!black},
    commentstyle=\color{green!50!black},
    stringstyle=\color{orange!80!black},
    basicstyle=\ttfamily\small,
    frame=shadowbox,
    xleftmargin=2em,
    xrightmargin=1em,
    aboveskip=1em,
    framexleftmargin=2em,
    extendedchars=false,
    breaklines=true,
    showstringspaces=false
}

\begin{document}

\begin{center}
{\huge \textbf{数值分析第一周上机作业}}

\vspace{0.3cm}
{\large 学号：241098117\quad，姓名：张城宗\quad}
\end{center}

\section{题目一：机器精度、下溢值与上溢值}

\subsection{问题}

分别针对 \textbf{双精度（double, float64）}和\textbf{单精度（single, float32）}浮点数，利用算法找出所使用计算机的机器精度 $\varepsilon_{\mathrm{mach}}$、下溢值与上溢值，并编程验证算法的正确性，比较两种精度的差异．

\subsection{算法思路}

IEEE 754 标准规定了两种常用浮点格式，结构如下：

\begin{table}[H]
\centering
\begin{tabular}{lccccc}
\toprule
格式 & 总位数 & 符号位 & 指数位 & 尾数位 & 指数偏置 \\
\midrule
双精度 (double) & 64 & 1 & 11 & 52 & 1023 \\
单精度 (single) & 32 & 1 & 8  & 23 & 127  \\
\bottomrule
\end{tabular}
\end{table}

\noindent\textbf{隐含位（Hidden Bit）．}规格化浮点数的形式为 $1.f \times 2^e$，整数部分的 $1$ 对所有规格化数恒成立，因此 IEEE 754 \textbf{不实际存储该位}，约定其隐式存在——称为\textbf{隐含位}．由此，虽然双精度只存储52位尾数，实际有效位数却是 $52+1=53$ 位，单精度为 $23+1=24$ 位．三个极限的理论值由有效位数 $p$（含隐含位）和指数范围 $[e_{\min}, e_{\max}]$ 给出：
$$
\varepsilon_{\mathrm{mach}} = 2^{-p}, \quad
x_{\min} = 2^{e_{\min}}, \quad
x_{\max} = (2-2^{-p})\cdot 2^{e_{\max}}.
$$

\noindent\textbf{浮点数的本质是二进制．}计算机内部存储的是64位（或32位）二进制串，并非十进制小数．MATLAB 显示的 \texttt{2.2204e-16} 只是将二进制值转换为十进制字符串以便人类阅读．可用 \texttt{num2hex(1.0)} 查看其十六进制存储值（\texttt{3FF0000000000000}），还原后即 $1.000\cdots0_2 \times 2^0$．"单/双精度"描述的是二进制表示的位宽，与是否用十进制显示无关．

代入具体参数，得到表 \ref{tab:theory} 中的理论值．三个算法的结构对双精度和单精度\textbf{完全相同}，仅数据类型不同；在 MATLAB 中，对单精度只需将初始值用 \texttt{single()} 转换，算术运算便自动在单精度下进行．

\subsubsection*{算法一：机器精度}

\textbf{思路：}$\varepsilon_{\mathrm{mach}}$ 的定义是满足 $1.0 + \varepsilon > 1.0$ 的最小正数．从 $\varepsilon=1$ 出发，不断折半，直到 $1+\varepsilon/2$ 在计算机中与 $1$ 相等，此时的 $\varepsilon$ 即为所求．

\begin{algorithm}[H]
    \caption{求机器精度}
    \begin{algorithmic}[1]
        \State $\varepsilon \e 1.0$
        \While{$\mathrm{fl}(1.0 + \varepsilon/2) > 1.0$}
            \State $\varepsilon \e \varepsilon / 2$
        \EndWhile
        \State \Return $\varepsilon$
    \end{algorithmic}
\end{algorithm}

\rmk IEEE 754 默认采用\textbf{最近舍入}（round-to-nearest-even）：若真实值恰好落在两个相邻浮点数的正中间，则舍入到尾数末位为偶数的那个．由此，任意实数 $x$ 的浮点化误差满足
$$\frac{|\mathrm{fl}(x)-x|}{|x|} \le u, \quad u = \frac{\varepsilon_{\mathrm{mach}}}{2} = 2^{-53} \text{（双精度）},$$
$u$ 称为\textbf{单位舍入误差}（unit roundoff）．我们的算法得到的 $\varepsilon_{\mathrm{mach}}=2^{-52}$ 是1与右邻浮点数的间距；正因为最近舍入将距离不超过 $\varepsilon_{\mathrm{mach}}/2$ 的数归并到1，实际误差上界比间距小一倍，即 $u=\varepsilon_{\mathrm{mach}}/2$．

\subsubsection*{算法二：下溢值（正规化边界检测）}

\textbf{思路：}下溢值即最小正规格化浮点数 \texttt{realmin}，是正规化得以维持的边界．正规化数 $x$ 的相对精度恰好为 $\varepsilon_{\mathrm{mach}}$，即 $x + x\varepsilon_{\mathrm{mach}} > x$；一旦 $x$ 变为非规格化数，$x\varepsilon_{\mathrm{mach}}$ 会因小于最小可表示正数而舍入为零，导致 $x + x\varepsilon_{\mathrm{mach}} = x$，相对精度"消失"．利用此性质检测正规化边界：

\begin{algorithm}[H]
    \caption{求下溢值（正规化边界检测）}
    \begin{algorithmic}[1]
        \State $x \e 1.0$
        \While{$\mathrm{fl}\!\left(\dfrac{x}{2} + \dfrac{x}{2}\cdot\varepsilon_{\mathrm{mach}}\right) > \dfrac{x}{2}$}
            \Comment{$x/2$ 仍是正规化数}
            \State $x \e x/2$
        \EndWhile
        \State \Return $x$ \Comment{$x = \mathtt{realmin} = 2^{e_{\min}}$}
    \end{algorithmic}
\end{algorithm}

\rmk 下溢值是正规化得以维持的临界点，而非计算机能表示的绝对最小正数；算法的终止条件正是正规化相对精度消失的瞬间．

\subsubsection*{算法三：上溢值（两阶段二分逼近）}

\textbf{思路：}$x_{\max}$ 的二进制表示为 $1.\underbrace{11\cdots1}_{52}\times 2^{1023}$，即 $2^{1023}+2^{1022}+\cdots+2^{971}$．算法分两个阶段：

\textbf{第一阶段}：从 $x=1$ 出发，不断翻倍，找到最大的2的整数次幂 $2^{1023}$（再翻倍则溢出）．

\textbf{第二阶段}：令步长 $s = x/2$，依次尝试将 $s$ 加入 $x$——若不溢出则接受，反之则放弃；每轮步长减半，直到 $s < 1$ 为止．这本质上是逐位填充尾数的52个二进制位1．

\begin{algorithm}[H]
    \caption{求上溢值（两阶段）}
    \begin{algorithmic}[1]
        \State $x \e 1.0$
        \While{$\mathrm{fl}(2x) \neq +\infty$}
            \State $x \e 2x$ \Comment{第一阶段：找到 $2^{1023}$}
        \EndWhile
        \State $s \e x/2$
        \While{$s \geq 1$}
            \If{$\mathrm{fl}(x+s) \neq +\infty$}
                \State $x \e x + s$ \Comment{第二阶段：接受当前位为1}
            \EndIf
            \State $s \e s/2$
        \EndWhile
        \State \Return $x$
    \end{algorithmic}
\end{algorithm}

\subsection{结果分析}

运行程序，将双精度与单精度的算法结果分别与理论值和 MATLAB 内置常量对比，如表 \ref{tab:theory} 所示．

\begin{table}[htbp]
\centering
\caption{双精度与单精度浮点数极限值对比}
\label{tab:theory}
\small  % 表格用小字号，防止溢出
\begin{tabular}{llccc}
\toprule
\textbf{量} & \textbf{精度} & \textbf{算法计算值} & \textbf{MATLAB 内置值} & \textbf{理论值} \\
\midrule
\multirow{2}{*}{机器精度}
  & 双精度 & $2.2204\times10^{-16}$ & $\varepsilon=2.2204\times10^{-16}$ & $2^{-52}$ \\
  & 单精度 & $1.1921\times10^{-7}$  & $\varepsilon=1.1921\times10^{-7}$  & $2^{-23}$ \\
\midrule
\multirow{2}{*}{\makecell[l]{下溢值\\（最小正规化数）}}
  & 双精度 & $2.2251\times10^{-308}$ & \texttt{realmin}$=2.2251\times10^{-308}$ & $2^{-1022}$ \\
  & 单精度 & $1.1755\times10^{-38}$  & \texttt{realmin}$=1.1755\times10^{-38}$  & $2^{-126}$ \\
\midrule
\multirow{2}{*}{上溢值}
  & 双精度 & $1.7977\times10^{308}$ & \texttt{realmax}$=1.7977\times10^{308}$ & $(2-2^{-52})\cdot2^{1023}$ \\
  & 单精度 & $3.4028\times10^{38}$  & \texttt{realmax}$=3.4028\times10^{38}$  & $(2-2^{-23})\cdot2^{127}$ \\
\bottomrule
\end{tabular}
\end{table}

\noindent \textbf{机器精度}：两种精度下算法结果均与内置值完全吻合，相对误差为0．单精度的机器精度约为 $1.19\times10^{-7}$，比双精度（$2.22\times10^{-16}$）大约9个数量级，这直接反映了两者尾数位数的差异（23位 vs 52位）．

\noindent \textbf{下溢值}：算法结果与 \texttt{realmin} 完全一致，相对误差为0．算法在 $x/2$ 失去相对精度 $\varepsilon_{\mathrm{mach}}$ 时停止，精确捕捉到正规化的边界（双精度 $2^{-1022}$，单精度 $2^{-126}$）．

\noindent \textbf{上溢值}：两阶段算法对双精度和单精度均给出与 \texttt{realmax} 完全一致的结果，相对误差为0．双精度的动态范围（$\sim10^{308}$）远大于单精度（$\sim10^{38}$），原因在于指数位数不同（11位 vs 8位）．

\subsection{结论}

三个算法均以"折半/翻倍+临界条件"为核心，对双精度和单精度\textbf{通用}，仅需切换数据类型（MATLAB 中用 \texttt{single()} 转换），结果与 IEEE 754 理论值完全吻合．

对比两种精度，关键差异如下：
\begin{itemize}
  \item \textbf{精度}：双精度机器精度 $\approx 2.2\times10^{-16}$，单精度 $\approx 1.2\times10^{-7}$，相差约9个数量级；
  \item \textbf{动态范围}：双精度可表示至 $\sim10^{\pm308}$，单精度至 $\sim10^{\pm38}$；
  \item \textbf{存储开销}：单精度占用内存仅为双精度的一半（32位 vs 64位）．
\end{itemize}

实际应用中，选择精度需权衡：对精度和动态范围要求高的科学计算一般用双精度；深度学习等对速度和内存敏感的场景则常用单精度甚至更低精度．

%% ============================================================
%%  题目二
%% ============================================================
\clearpage

\section{题目二：避免相减相消计算 \texorpdfstring{$y = x - \sin x$}{y = x - sin(x)}}

\subsection{问题}

设计算法计算 $y = x - \sin x$，使得有效位的丢失最多 $1$ 位，并编程验证所设计算法的正确性．

\subsection{算法思路}

\noindent\textbf{相减相消（Cancellation）现象．}
当 $x$ 接近 $0$ 时，$\sin x \approx x$，两者几乎相等，直接相减会引发\textbf{相减相消}：
两个相近的数之差虽然绝对误差不大，但结果的相对误差可极大，导致大量有效位丢失．

由讲义\textbf{精度丢失定理（定理 1.3.1）}：设 $x, y > 0$ 为规格化二进制浮点数且 $x > y$，若
$$2^{-q} \le 1 - \frac{y}{x} \le 2^{-p},$$
则做减法 $x - y$ 时至少丢失 $p$ 个、至多丢失 $q$ 个二进制精度位．

\noindent\textbf{误差估计．}
对 $y = x - \sin x$，以 $x$ 为大数、$\sin x$ 为小数，利用 $\sin x$ 的 Taylor 展开得
$$1 - \frac{\sin x}{x} = \frac{x^2}{6} - \frac{x^4}{120} + \cdots \approx \frac{x^2}{6} \quad (x \to 0).$$
因此直接计算约丢失 $\log_2\!\left(\dfrac{6}{x^2}\right)$ 个二进制位．

\noindent\textbf{阈值的确定．}
要求直接公式至多丢失 $1$ 位，需
$$1 - \frac{\sin x}{x} \;\ge\; 2^{-1} = 0.5 \;\Longleftrightarrow\; \frac{\sin x}{x} \le \frac{1}{2}.$$
数值求解方程 $\sin x_0 / x_0 = 1/2$，得 $x_0 \approx 1.895$．
当 $|x| \ge x_0$ 时，直接公式至多丢失 $1$ 位；当 $|x| < x_0$ 时，需使用改进算法．

\noindent\textbf{改进算法：Taylor 级数．}
将 $\sin x$ 展开并代入：
$$x - \sin x = \frac{x^3}{3!} - \frac{x^5}{5!} + \frac{x^7}{7!} - \cdots
= \sum_{k=1}^{\infty} \frac{(-1)^{k+1}\, x^{2k+1}}{(2k+1)!}.$$
此级数各项均为有效加法，不存在相减相消，对小 $x$ 收敛极快（首项 $x^3/6$ 即为主项）．相邻两项的递推关系为
$$t_{k+1} = t_k \cdot \frac{-x^2}{(2k+2)(2k+3)},$$
逐项累加至 $|t_k| \le |y|\cdot\varepsilon_{\mathrm{mach}}$ 时停止，即可达到机器精度．

综合两种情形，改进算法如下：

\begin{algorithm}[H]
    \caption{无相减相消地计算 $y = x - \sin x$}
    \begin{algorithmic}[1]
        \Require $x$，切换阈值 $x_0 \approx 1.895$
        \Ensure $y = x - \sin x$（至多丢失 $1$ 个二进制位）
        \If{$|x| < x_0$}
            \State $t \e x^3 / 6,\quad y \e t,\quad k \e 1$
            \While{$|t| > |y| \cdot \varepsilon_{\mathrm{mach}}$}
                \State $t \e t \cdot (-x^2) / [(2k{+}2)(2k{+}3)],\quad y \e y + t,\quad k \e k{+}1$
            \EndWhile
        \Else
            \State $y \e x - \sin x$ \Comment{$|x| \ge x_0$，至多丢失 $1$ 位}
        \EndIf
        \State \Return $y$
    \end{algorithmic}
\end{algorithm}

\subsection{结果分析}

\noindent\textbf{直接公式的精度损失．}
以 Taylor 级数结果为参考真值，表 \ref{tab:p2} 展示了直接公式 $y = x - \sin x$ 在不同 $x$ 下的相对误差．

\begin{table}[H]
\centering
\caption{直接公式与 Taylor 级数的对比（相对误差 $\approx 12\varepsilon_{\mathrm{mach}}/x^2$）}
\label{tab:p2}
\small
\begin{tabular}{cccc}
\toprule
$x$ & 直接公式 & Taylor 级数（精确） & 相对误差 \\
\midrule
$10^{-2}$ & $1.66667\times10^{-7}$ & $1.66667\times10^{-7}$ & $\sim 10^{-12}$ \\
$10^{-4}$ & $1.66667\times10^{-13}$ & $1.66667\times10^{-13}$ & $\sim 10^{-8}$ \\
$10^{-6}$ & $1.66533\times10^{-19}$ & $1.66667\times10^{-19}$ & $\sim 10^{-3}$ \\
$10^{-8}$ & $0$ & $1.66667\times10^{-25}$ & $1$ （完全失真）\\
\bottomrule
\end{tabular}
\end{table}

相对误差的理论近似为 $12\varepsilon_{\mathrm{mach}}/x^2$，在双对数坐标下呈斜率 $-2$ 的直线，
阈值线 $x = x_0 \approx 1.895$ 恰与误差曲线交于 $y \approx 3\varepsilon_{\mathrm{mach}}$，
即「至多丢失 $1$ 个二进制位」的临界点；阈值右侧误差已降至机器精度量级，直接公式完全可用．

\subsection{结论}

直接计算 $x - \sin x$ 对小 $|x|$ 会产生严重的相减相消，丢失约 $\log_2(6/x^2)$ 个有效二进制位．改进算法以 Taylor 级数替代直接公式，彻底消除相消现象，在全域精度均达到机器精度量级．切换阈值 $x_0 \approx 1.895$ 由精度丢失定理严格导出，保证了改进算法至多丢失 $1$ 个二进制位．

%% ============================================================
%%  题目三
%% ============================================================
\clearpage

\section{题目三：积分递推公式的数值不稳定性}

\subsection{问题}

计算积分序列
$$y_n = \int_0^1 x^n e^x\,\mathrm{d}x, \quad n \ge 0.$$
由分部积分可得递推公式
$$y_{n+1} = \mathrm{e} - (n+1)\,y_n.$$
编程验证按上述公式设计的正向递推算法不稳定，并给出稳定的替代算法．

\subsection{算法思路}

\noindent\textbf{初始值与真值界．}
$$y_0 = \int_0^1 e^x\,\mathrm{d}x = \mathrm{e} - 1 \approx 1.71828.$$
由于 $x \in [0,1]$ 时 $1 \le e^x \le \mathrm{e}$，有
$$\frac{1}{n+1} \le y_n \le \frac{\mathrm{e}}{n+1},$$
即 $y_n > 0$ 且单调趋近于 $0$．

\noindent\textbf{正向递推的误差分析（不稳定）．}
设 $\varepsilon_n = \tilde{y}_n - y_n$ 为第 $n$ 步的累积误差．由递推公式：
$$\varepsilon_{n+1} = -(n+1)\,\varepsilon_n \;\Longrightarrow\; |\varepsilon_n| = n!\,|\varepsilon_0|.$$
即使初始误差 $|\varepsilon_0|$ 仅为机器精度量级（$\sim 2\times10^{-16}$），经 $n$ 步后误差放大为 $n!\,|\varepsilon_0|$．当 $n = 20$ 时，$20! \approx 2.4 \times 10^{18}$，误差已完全淹没真值——算法\textbf{不稳定}．

\noindent\textbf{反向递推（稳定）．}
由递推公式解出 $y_n$：
$$y_n = \frac{\mathrm{e} - y_{n+1}}{n+1}.$$
误差传播：$\varepsilon_n = -\varepsilon_{n+1}/(n+1)$，即每步误差\textbf{缩小} $n+1$ 倍．
取足够大的 $N$（如 $N=40$），令 $\tilde{y}_N = 0$（起点误差 $\le \mathrm{e}/(N+1) \approx 0.066$），
反向递推到 $n=0$ 时，误差缩小为 $|\varepsilon_0| \approx |\varepsilon_N|/N! \approx 0$——算法\textbf{稳定}．

两种算法的误差传播对比如下：

\begin{table}[H]
\centering
\caption{正向与反向递推的误差传播对比}
\begin{tabular}{lll}
\toprule
 & 正向递推 & 反向递推 \\
\midrule
递推方向 & $n=0 \to N$ & $n=N \to 0$ \\
误差传播 & $|\varepsilon_{n+1}| = (n+1)|\varepsilon_n|$ & $|\varepsilon_n| = |\varepsilon_{n+1}|/(n+1)$ \\
增长规律 & $|\varepsilon_n| = n!\,|\varepsilon_0|$（指数增长） & $|\varepsilon_n| \to 0$（指数衰减） \\
稳定性 & \textbf{不稳定} & \textbf{稳定} \\
\bottomrule
\end{tabular}
\end{table}

\subsection{结果分析}

表 \ref{tab:p3} 展示了两种递推算法与 MATLAB \texttt{integral()} 参考真值的对比：

\begin{table}[H]
\centering
\caption{积分 $y_n$ 的正向/反向递推与参考真值对比}
\label{tab:p3}
\small
\begin{tabular}{cccc}
\toprule
$n$ & 参考真值 & 正向递推计算值 & 反向递推计算值 \\
\midrule
0  & $1.718281828$ & $1.718281828$ & $1.718281828$ \\
5  & $0.024313660$ & $0.024313660$ & $0.024313660$ \\
10 & $0.016105880$ & $0.016108\cdots$（开始偏离） & $0.016105880$ \\
15 & $0.013103970$ & $-12.17\cdots$（已完全错误） & $0.013103970$ \\
20 & $0.011109440$ & $2.7\times10^{6}$（溢出量级） & $0.011109440$ \\
\bottomrule
\end{tabular}
\end{table}

正向递推在 $n \approx 12$ 时首次出现负值（真值始终为正），之后完全失真；反向递推在整个区间内与参考真值完全吻合，验证了稳定性分析的结论．

\subsection{结论}

正向递推公式 $y_{n+1} = \mathrm{e} - (n+1)y_n$ 数值不稳定，误差按 $n!$ 增长，实用中应\textbf{禁止使用}．
改用反向递推 $y_n = (\mathrm{e} - y_{n+1})/(n+1)$，从大 $n$ 处的近似起点出发，误差逐步缩小，可得到全程精确的结果．
这一例子表明：递推公式的数学等价性不等于数值等价性，方向的选择决定了算法的稳定性．

%% ============================================================
%%  题目四
%% ============================================================
\clearpage

\section{题目四：实数序列递推算法的稳定性}

\subsection{问题}

考察由以下递推关系定义的实数序列 $\{x_n\}$：
$$x_{n+1} = \frac{13}{3}\,x_n - \frac{4}{3}\,x_{n-1}, \quad n = 1, 2, 3, \ldots$$
分别在两组初值下分析算法的数值稳定性：
\begin{itemize}
  \item \textbf{情形一：}$x_0 = 1,\ x_1 = 1/3$
  \item \textbf{情形二：}$x_0 = 1,\ x_1 = 4$
\end{itemize}

\subsection{算法思路}

\noindent\textbf{特征方程与通解．}
令 $x_n = r^n$ 代入递推关系，除以 $r^{n-1}$（$r \ne 0$），得特征方程
$$3r^2 - 13r + 4 = 0 \;\Longrightarrow\; r = \frac{13 \pm 11}{6},$$
故两个特征根为 $r_1 = 4,\ r_2 = 1/3$．通解为
$$x_n = A \cdot 4^n + B \cdot \left(\frac{1}{3}\right)^n.$$

\noindent\textbf{情形一的精确解与稳定性分析．}
代入初值 $x_0 = 1,\ x_1 = 1/3$，解得 $A = 0,\ B = 1$，故精确解
$$x_n = \left(\frac{1}{3}\right)^n \xrightarrow{n\to\infty} 0.$$
然而，浮点计算中 $x_1 = 1/3$ 存在舍入误差 $\delta \sim \varepsilon_{\mathrm{mach}}$，等价于引入一个微小的 $r_1 = 4$ 分量扰动，使实际计算值为
$$\tilde{x}_n = \left(\frac{1}{3}\right)^n + \delta \cdot 4^n.$$
由于 $4^n$ 指数增长，而真值 $(1/3)^n$ 指数衰减，相对误差
$$\frac{|\tilde{x}_n - x_n|}{|x_n|} = \delta \cdot 12^n \xrightarrow{n\to\infty} \infty.$$
算法\textbf{不稳定}（$r_1 = 4$ 为增长的"寄生解"，被微小扰动激发）．

\noindent\textbf{情形二的精确解与稳定性分析．}
代入初值 $x_0 = 1,\ x_1 = 4$，解得 $A = 1,\ B = 0$，故精确解
$$x_n = 4^n.$$
浮点扰动激发 $r_2 = 1/3$ 分量，使计算值为 $\tilde{x}_n = 4^n + \delta\cdot(1/3)^n$，相对误差
$$\frac{|\tilde{x}_n - x_n|}{|x_n|} = \delta\cdot\left(\frac{1}{12}\right)^n \xrightarrow{n\to\infty} 0.$$
算法\textbf{稳定}（扰动项 $(1/3)^n$ 相对于主项 $4^n$ 迅速衰减）．

两种情形的对比如下：

\begin{table}[H]
\centering
\caption{两种初值下的稳定性分析}
\begin{tabular}{lllll}
\toprule
情形 & 初值 & 精确解 & 主特征根 & 相对误差增长规律 \\
\midrule
一 & $x_0=1,x_1=1/3$ & $(1/3)^n \to 0$ & $r_2 = 1/3$（衰减） & $\delta \cdot 12^n \to \infty$（\textbf{不稳定}）\\
二 & $x_0=1,x_1=4$   & $4^n \to \infty$ & $r_1 = 4$（增长）   & $\delta \cdot (1/12)^n \to 0$（\textbf{稳定}）\\
\bottomrule
\end{tabular}
\end{table}

\subsection{结果分析}

表 \ref{tab:p4} 展示了两种情形的数值计算结果与精确解的比较：

\begin{table}[H]
\centering
\caption{两种初值下的数值计算值、精确值与相对误差}
\label{tab:p4}
\small
\begin{tabular}{ccccc}
\toprule
\multirow{2}{*}{$n$} & \multicolumn{2}{c}{情形一：$x_n = (1/3)^n$（应趋于0）} & \multicolumn{2}{c}{情形二：$x_n = 4^n$（应趋于$\infty$）} \\
\cmidrule(lr){2-3}\cmidrule(lr){4-5}
 & 递推计算值 & 相对误差 & 递推计算值 & 相对误差 \\
\midrule
5  & $4.115\times10^{-3}$ & $\sim 10^{-11}$ & $1024$ & $\sim 10^{-13}$ \\
10 & $1.694\times10^{-5}$ & $\sim 10^{-5}$ & $1048576$ & $\sim 10^{-15}$ \\
20 & $-4.8\times10^{-3}$（负值！）& $\sim 10^{5}$ & $1.10\times10^{12}$ & $\sim 10^{-15}$ \\
30 & $1.7\times10^{4}$（完全发散）& $\gg 1$ & $1.15\times10^{18}$ & $\sim 10^{-14}$ \\
\bottomrule
\end{tabular}
\end{table}

情形一在 $n \approx 18$ 时首次出现负值（真值恒为正），此后完全发散；情形二在整个计算过程中相对误差维持在机器精度量级，算法完全稳定．

\subsection{结论}

同一递推公式在不同初值下可以表现出截然不同的数值稳定性，其根本原因在于特征根的大小关系：
\begin{itemize}
  \item 若精确解对应\textbf{较小的特征根}（如 $|r_2|=1/3$），任何浮点扰动都会激发较大特征根（$|r_1|=4$）对应的"寄生解" $4^n$，导致误差以 $12^n$ 速率爆炸增长——\textbf{算法不稳定}；
  \item 若精确解对应\textbf{较大的特征根}（如 $|r_1|=4$），扰动激发的小特征根分量 $(1/3)^n$ 相对于主项迅速衰减——\textbf{算法稳定}．
\end{itemize}

一般规律：对于线性递推方程，若算法要求追踪的解分量对应\textbf{非主导特征根}（绝对值较小的根），则数值误差将被主导特征根的增长所放大，产生不稳定性；只有当追踪的解分量本身是\textbf{主导分量}时，算法才具有数值稳定性．

%% ============================================================
%%  附录
%% ============================================================
\clearpage

\section{附录：程序代码}

\subsection{题目一}

\begin{lstlisting}[style=matlab, caption={problem1\_machine\_precision.m（核心部分）}]
clear; clc;

%% --- 双精度 ---
% 机器精度
eps_calc = 1.0;
while (1.0 + eps_calc / 2) > 1.0
    eps_calc = eps_calc / 2;
end

% 下溢值：正规化边界检测（x/2 失去相对精度 eps 时停止）
underflow_calc = 1.0;
while underflow_calc/2 + (underflow_calc/2)*eps > underflow_calc/2
    underflow_calc = underflow_calc / 2;
end

% 上溢值（两阶段）
overflow_calc = 1.0;
while ~isinf(overflow_calc * 2),  overflow_calc = overflow_calc * 2;  end
step = overflow_calc / 2;
while step >= 1
    if ~isinf(overflow_calc + step),  overflow_calc = overflow_calc + step;  end
    step = step / 2;
end

fprintf('[双精度] eps=%.4e, underflow=%.4e, overflow=%.4e\n', ...
    eps_calc, underflow_calc, overflow_calc);

%% --- 单精度：仅需把初始值换为 single() ---
eps_s = single(1.0);
while single(1.0) + eps_s / single(2.0) > single(1.0)
    eps_s = eps_s / single(2.0);
end

underflow_s = single(1.0);
while underflow_s/single(2) + (underflow_s/single(2))*eps_s > underflow_s/single(2)
    underflow_s = underflow_s / single(2);
end

overflow_s = single(1.0);
while ~isinf(overflow_s * single(2.0)),  overflow_s = overflow_s * single(2.0);  end
step_s = overflow_s / single(2.0);
while step_s >= single(1.0)
    if ~isinf(overflow_s + step_s),  overflow_s = overflow_s + step_s;  end
    step_s = step_s / single(2.0);
end

fprintf('[单精度] eps=%.4e, underflow=%.4e, overflow=%.4e\n', ...
    eps_s, underflow_s, overflow_s);
\end{lstlisting}

\subsection{题目二核心代码}

\begin{lstlisting}[style=matlab, caption={problem2\_avoid\_cancellation.m（核心部分）}]
THRESHOLD = 1.895;  % 阈值：sin(x)/x = 0.5 的数值解

% 改进算法：根据 |x| 与阈值自动切换公式
function y = improved_formula(x, threshold)
    if abs(x) < threshold
        y = taylor_series(x);   % 小 x：Taylor 级数，无相减相消
    else
        y = x - sin(x);         % 大 x：直接公式，至多丢失 1 位
    end
end

% Taylor 级数：x - sin(x) = x^3/6 - x^5/120 + ...
function y = taylor_series(x)
    if x == 0;  y = 0;  return;  end
    term = x^3 / 6;    % 首项 k=1
    y = term;
    k = 1;
    while abs(term) > abs(y) * eps
        term = term * (-x^2) / ((2*k+2) * (2*k+3));  % 递推到第 k+1 项
        y = y + term;
        k = k + 1;
        if k > 200; break; end
    end
end
\end{lstlisting}

\subsection{题目三核心代码}

\begin{lstlisting}[style=matlab, caption={problem3\_integral\_recursion.m（核心部分）}]
e_val = exp(1);
N_forward  = 25;   % 正向递推阶数
N_backward = 40;   % 反向递推起始阶数

%% 参考真值（MATLAB 自适应数值积分）
for n = 0:N_forward
    y_true(n+1) = integral(@(x) x.^n .* exp(x), 0, 1, 'RelTol', 1e-15);
end

%% 正向递推：y_{n+1} = e - (n+1)*y_n（不稳定）
y_forward(1) = e_val - 1;     % y_0 = e - 1
for n = 0 : N_forward - 1
    y_forward(n+2) = e_val - (n+1) * y_forward(n+1);
end

%% 反向递推：y_n = (e - y_{n+1}) / (n+1)（稳定）
y_backward(N_backward + 1) = 0;   % 起点近似：y_N ~= 0
for n = N_backward - 1 : -1 : 0
    y_backward(n+1) = (e_val - y_backward(n+2)) / (n+1);
end
\end{lstlisting}

\subsection{题目四核心代码}

\begin{lstlisting}[style=matlab, caption={problem4\_sequence\_stability.m（核心部分）}]
c1 = 13/3;   c2 = 4/3;   N = 60;

%% 情形一：x0=1, x1=1/3，精确解 x_n = (1/3)^n
x_c1(1) = 1;   x_c1(2) = 1/3;
for n = 2:N
    x_c1(n+1) = c1 * x_c1(n) - c2 * x_c1(n-1);
end
x_exact_c1 = (1/3).^(0:N);   % 精确解

%% 情形二：x0=1, x1=4，精确解 x_n = 4^n
x_c2(1) = 1;   x_c2(2) = 4;
for n = 2:N
    x_c2(n+1) = c1 * x_c2(n) - c2 * x_c2(n-1);
end
x_exact_c2 = 4.^(0:N);       % 精确解
\end{lstlisting}

\end{document}
